% !TeX program = xelatex
\documentclass[11pt, a4paper]{article}
\usepackage[utf8]{inputenc}
\usepackage[T1]{fontenc}
\usepackage{xeCJK} % For Chinese support if needed, but the review is in English as per academic standard
\usepackage[margin=2.5cm]{geometry}
\usepackage{xcolor}
\usepackage{titlesec}
\usepackage{enumitem}
\usepackage{tabularx}
\usepackage{booktabs}
\usepackage{amsmath, amssymb, amsthm}
\usepackage{hyperref}

% Define custom colors for review categories
\definecolor{NavyBlue}{HTML}{1E3A8A}
\definecolor{RustOrange}{HTML}{9A3412}
\definecolor{ForestGreen}{HTML}{065F46}
\definecolor{DeepRed}{HTML}{991B1B}

% Setup hyperefs
\hypersetup{
    colorlinks=true,
    linkcolor=NavyBlue,
    filecolor=magenta,      
    urlcolor=NavyBlue,
    pdftitle={Review and Correction Documentation},
}

% Section styling
\titleformat{\section}{\Large\bfseries\color{NavyBlue}}{}{0em}{}[\titlerule]
\titleformat{\subsection}{\large\bfseries\color{NavyBlue!80}}{}{0em}{}

\begin{document}

% --- Title Page ---

\begin{titlepage}
    \vspace*{\fill}
    \begin{center}
        % 保持主标题巨大且加粗
        {\Huge \bfseries Review and Correction Documentation \par}
        \vspace{2.5cm} % 增加标题与下方信息块的间距，防止拥挤

        % 将字号从 \large 改为 \Large (甚至 \LARGE)
        {\Large 
        \renewcommand{\arraystretch}{1.8} % 增大行高以匹配大字体
        \begin{tabular}{r l}
            \textbf{Course:} & Academic Writing and Communication in English \\
            \textbf{Reviewer:} & Nie Jiaduo, Wen Zijie, and Zhong Haoyu \\
            \textbf{Institution:} & Department of Mathematics, SUSTech \\
        \end{tabular}
        \par}
    \end{center}
    \vspace*{\fill}
\end{titlepage}

\newpage
\tableofcontents
\newpage






% --- Section 1: Classification System ---
\section{Review Framework and Overview}

% --- 1.1 目标论文基础信息 ---
\subsection{Target Paper Information}

The core metadata of the manuscript selected for this critical review is summarized in the table below:

\begin{center}
    \renewcommand{\arraystretch}{1.5} % 保持行距美观
    \begin{tabularx}{0.95\textwidth}{r X} % 使用 X 列实现长标题自动换行对齐
        \textbf{Title:} & Projection, Measure, and Idempotent Relations: Independent Axioms and a Fixed-Point Coupling Law \\
        \textbf{Field:} & Mathematical Logic \\
        \textbf{Author:} & Yunbeom Yi \\
        \textbf{arXiv ID:} & \href{https://arxiv.org/abs/2604.18640}{2604.18640} \\
        \textbf{Submission Date:} & April 19, 2026 \\
        \textbf{Background:} & Category Theory, Measure Theory, Set Theory, and Operator Theory \\
    \end{tabularx}
\end{center}

\vspace{0.5cm} % 与下一小节的间距

\subsection{Classification of Errors}


To ensure a structured and comprehensive review, we categorize the identified errors into four distinct types based on their nature and severity.

\begin{table}[htbp]
\centering
\renewcommand{\arraystretch}{1.5}
\begin{tabularx}{\textwidth}{l >{\raggedright\arraybackslash}p{3.5cm} X}
\toprule
\textbf{Marker Color} & \textbf{Category} & \textbf{Definition} \\
\midrule
\textcolor{NavyBlue}{\textbf{Navy Blue}} & \textbf{Language \& Punctuation} & \textbf{Violations of standard grammatical rules and typographical conventions.} Encompasses basic syntactic deficiencies, typos, and improper spacing. \\

\textcolor{RustOrange}{\textbf{Rust Orange}} & \textbf{Style \& Architecture} & \textbf{Impediments to narrative flow and structural clarity of academic expression.} Addresses disorganized exposition, abrupt logical leaps, and keyword overloading. \\

\textcolor{ForestGreen}{\textbf{Forest Green}} & \textbf{Math Formatting} & \textbf{Deviations from rigorous mathematical typesetting standards.} Identifies missing punctuation in equations and improper \LaTeX~notation. \\

\textcolor{DeepRed}{\textbf{Deep Red}} & \textbf{Logical Flaws} & \textbf{Fundamental breaches of mathematical reasoning and conceptual integrity.} Encompasses axiomatic redundancy and circular reasoning. \\
\bottomrule
\end{tabularx}
\end{table}
Throughout the annotated PDF file of the manuscript, these corresponding colors are utilized to highlight and distinguish the specific types of errors.

\newpage 
% --- Section 2: Detailed Review ---
\section{Comprehensive Review, Analysis, and Extension}

\subsection{Title Analysis}


The title of the reviewed manuscript is presented below:

\begin{center}
    \textit{"Projection, Measure, and Idempotent Relations: Independent Axioms and a Fixed-Point Coupling Law"}
\end{center}

\begin{itemize}
    \item \textcolor{RustOrange}{\textbf{[Style \& Architecture] Noun Stacking:}} The current title is essentially a \textbf{stacking of nouns} that lacks a clear logical hierarchy. This "noun-piling" structure is purely descriptive and \textbf{fails to reflect the core work} or the central contribution of the research.
    
    \item \textbf{Correction Suggestion:} The title should be restructured to highlight the primary scientific result. A more effective version would be: \textit{"A Fixed-Point Coupling Law for Idempotent Relations in Measure Theory"}.
\end{itemize}
The title is the first element encountered by a reader and should provide a clear focus on the paper's core contribution.
In a word, \textbf{the title should convey some information to the reader.}

The following are the requirements for titles according to APA Style:
\begin{center}
    \small \textit{``There is no maximum length for titles; however, keep titles focused and include key terms.''} \\
    --- \textbf{APA Style (7th Edition)}
\end{center}
\begin{center}
    \small \textit{``Keep it shorter than 12 words''} \\
    --- \textbf{APA Style (6th Edition)}
\end{center}



\subsection{Abstract Analysis}
The abstract serves as the gateway to the paper. It should provide a high-level summary of the motivation, methodology, and main results. However, the current abstract suffers from several significant structural and typographical issues:

\begin{itemize}
    \item \textcolor{RustOrange}{\textbf{[Style \& Architecture] Technical Overload \& Lack of Big Picture:}} The abstract dives immediately into dense technical details and a minimal axiom system without providing any global context or motivation. It completely lacks a general introduction to the overall work, making it highly inaccessible to a broader mathematical audience.

    \item \textcolor{RustOrange}{\textbf{[Style \& Architecture] Formula Clutter \& Undefined Notation:}} The text is excessively cluttered with mathematical formulas and complex tuples, such as \[(X, \mathcal{A}, \mu, \mu^{\otimes 2}, R, I, \Pi_R, G, E_0, \eta).\] More critically, several symbols within this tuple (e.g., $E_0$) are presented without any explanation or definition, which severely breaks the self-contained nature required of an abstract.

    \item \textcolor{NavyBlue}{\textbf{[Language \& Punctuation] Hyphenation \& Typographical Errors:}} As highlighted in the annotated manuscript, there is a pervasive misuse of hyphens in mathematical terminology (e.g., \textit{``ZFC-internal''}, \textit{``projection-measure''}, \textit{``Banach-contraction''}). Additionally, there are basic capitalization errors, such as using lowercase letters immediately following a colon (e.g., \textit{``Consistency: the''} should be \textit{``Consistency: The''}).
    
    \item \textcolor{RustOrange}{\textbf{[Style \& Architecture] Keyword Redundancy:}} The list of keywords is excessively long and redundant. Keywords are intended to efficiently index the paper's primary contributions for search engines and databases. An overloaded and repetitive list dilutes the core focus and diminishes its utility.

    \item \textbf{Correction Suggestion:} The abstract must be extensively rewritten. It should begin with 1-2 sentences establishing the general background and the problem being solved. Complex tuples and undefined symbols should be removed in favor of descriptive textual summaries. Finally, all hyphenation and capitalization errors must be systematically corrected and the keyword list should be trimmed to only the most essential terms.
\end{itemize}


The manuscript provides ten keywords:

\begin{center}
    \textit{Set Theory (ZFC); Measure Theory; Charges and Finitely Additive Measures; Idempotent Relations; Retraction Operators; Coupling Law; Banach Fixed-Point Theorem; Independence of Axioms; Categorical Structure; Quotient Factorization.} 
\end{center}



Refer to standard academic formatting guidelines regarding keyword selection:

\begin{center}
    \small \textit{``Keywords need to be descriptive and capture the most important aspects of your paper. They are used for indexing in databases and as search terms for readers. \textbf{Include three to five words, phrases, or acronyms as keywords.}''} \\
    --- \textbf{APA Style(7th Edition)}
\end{center}

The current list exceeds the 3-5 limit and is highly redundant. Conceptually, ``Measure Theory'' and ``Charges and Finitely Additive Measures'' overlap. ``Set Theory'' is too broad for a paper focused strictly on ZFC, and ``Categorical Structure'' is too vague to be informative.

Trim the list to a maximum of five specific terms. A suggested revision: 
\begin{center}
\textit{Finitely Additive Measures; Idempotent Relations; Fixed-Point Coupling Law; Axiom Independence; ZFC}.
\end{center}

\subsection{Introduction Analysis}
The introduction should provide a clear and professionally formatted overview of the paper's framework. However, an analysis of this section reveals several recurring stylistic and typographical errors:

\begin{itemize}
    \item \textcolor{RustOrange}{\textbf{[Stylistic \& Structural] Terminology Inconsistency:}} The manuscript refers to the structure as a ``set algebra'' in the introduction , but later formally defines it as an ``algebra of sets'' (Definition 2.1). Mathematical terminology must remain strictly uniform throughout the text to avoid ambiguity.
    
    \item \textcolor{NavyBlue}{\textbf{[Language \& Punctuation] Improper Use of Slashes:}} The author frequently uses slashes to combine concepts, such as ``projection/retraction operators''. In formal academic writing, slashes should be avoided in prose. They must be replaced with precise conjunctions (e.g., ``projection and retraction'').
    
    \item \textcolor{NavyBlue}{\textbf{[Language \& Punctuation] Missing Hyphens in Compound Modifiers:}} Compound adjectives preceding a noun must be hyphenated to prevent misreading. For instance, the phrase ``measure preserving retraction''  lacks a hyphen and should be corrected to ``measure-preserving retraction''.
    \item \textcolor{RustOrange}{\textbf{[Style \& Architecture] Abrupt Opening \& Structural Inversion:}} The introduction opens abruptly by presenting a dense 10-tuple axiom system in the very first sentence. Standard academic structure requires establishing the general background and research motivation before diving into technical formalities. The actual motivation is awkwardly delayed until Section 1.1.
    
    \item \textbf{Correction Suggestion:} Reorder the introductory content. Begin the section with the context currently found in Section 1.1 (Motivation) to properly orient the reader, and only then introduce the formal mathematical framework.
\end{itemize}
An answer on MathOverflow\footnote{\url{https://mathoverflow.net/questions/254372/arrangement-of-the-introduction-of-a-paper}} explicitly outlines four core principles for arranging a paper's introduction:
\begin{itemize}
    \item the results are easy to grasp;
    \item the results are clearly motivated;
    \item the reader gets convinced of the significance of the results;
    \item the reader gets interested to read the rest of the paper.
\end{itemize}

In practice, the manuscript fails to satisfy any of these four principles. 

First, the introduction is inaccessible due to excessive technical density; the immediate presentation of a complex 10-tuple without prior context makes the results \textbf{difficult for the reader to grasp.}

Regarding the second and third principles, the provided motivation is superficial and unconvincing. Although the author claims connections to various mathematical domains—including measure theory, topology, and analysis—these relationships are merely listed rather than being synthetically elaborated. Consequently, \textbf{the significance of the work is not convincingly established, leaving the reader skeptical of its actual impact}.

Finally, the introduction fails to engage the reader because the \textbf{"Contributions" section lacks a clear bridge to classical cases.} Specifically, the author does not clarify the conditions under which the two critical hypotheses—fiber measurability and regularity—are satisfied. Since the theory's overall utility hinges directly on the validity of these assumptions, this lack of clarity discourages further investigation.

Clark Barwick's {advice}\footnote{\url{https://bimsa.net/activity/MathematicalCommunication/files/Clark\%20Barwick\%20--\%20Notes\%20on\%20mathematical\%20writing.pdf}} on writing an introduction is also highly useful. He specifically divides the introduction into the following steps:

\begin{quote}
    \small
    \begin{enumerate}
        \item \textbf{Establish common ground:} Tell a story your target audience knows (or should know) very well.
        \item \textbf{Identify the gap:} Highlight aspects of this story that are mysterious, incomplete, or inadequately explained.
        \item \textbf{State contributions explicitly:} State your contributions as Theorems with cross-references to their proofs.
        \item \textbf{Connect the dots:} Explain how your contributions address the gaps raised in step 2.
        \item \textbf{Iterate:} Repeat steps 1--4 if needed.
        \item \textbf{Stop:} Get on with the paper already.
    \end{enumerate}
\end{quote}



\subsection{Conclusion Analysis}

The manuscript's conclusion is structured as follows:

\begin{quote}
    \small \textit{
    We have set up a minimal ZFC-internal axiom system for admissible structural models pre-structural data $(X, \mathcal{A}, \mu, \mu^{\otimes 2}, R, I, \Pi_R, G, E_0, \eta)$ satisfying Axioms I, II, and III and studied its basic properties. \\
    \indent The axiom system is consistent in ZFC (Section 5), and its three axioms, together with the three subclauses of Axiom III, are mutually independent (Theorem 6.7, Corollary 6.8). The coupling law of Axiom III admits an operator-theoretic reformulation as the unique bounded finitely additive solution of the Banach-contraction equation $f = T_\eta f$, with closed form $f_*(B) = \mu(B) + \frac{\eta}{1-\eta}\mu(\Pi_R^{-1}(B))$ (Theorem 4.8). Admissible structural models with a common $\eta$ form a category $\mathbf{Struct}_\eta$ (Section 8); under fiber measurability and either (R-fin) or (R-ctbl), the general admissibility problem reduces to the identity-retraction case (Theorem 8.18), where the classification of Theorem 5.4 yields the block dichotomy $\mu(C_k) \in \{0, (1-\eta)^{-1}\}$. Under $\sigma$-additive probability normalization, the $B = X$ instance of the coupling law forces $\eta = 0$ (Section 7), so $\eta \neq 0$ families live in the finitely additive setting or among non-probability $\sigma$-additive measures. \\
    \indent \textbf{Further directions.} Several extensions suggest themselves: replacing $\mathcal{A}$ by a $\sigma$-algebra under hypotheses giving a well-defined product charge; enriching $G$ to a weighted kernel $K : X \times X \to [0, \infty)$ and comparing with graph limits [15]; further categorical development, in particular an adjoint characterization of the restriction $\rho : \mathbf{Struct}_\eta^{\text{fib}} \to \mathbf{Struct}_\eta^{\text{id}}$ [17, 20]; operator-theoretic variants replacing $\mu$ by a state or trace; and scaling limits $X_n \to X$ controlled by local instances of the coupling law. Removing the fiber hypotheses in the quotient-factorization reduction, for instance in the uncountable $\sigma$-additive case, also requires further assumptions such as separability of $\mathcal{A}$.
    }
\end{quote}


The conclusion of a paper should 
\begin{itemize}
    \item elevate the reader's perspective,
    \item synthesizing the core insights,
    \item and discussing their broader implications.
\end{itemize}

The current conclusion, however, fails to achieve this and suffers from significant structural limitations:

\begin{itemize}
    \item \textcolor{RustOrange}{\textbf{[Style \& Architecture] Redundant Summary \& Lack of Macroscopic Perspective:}} The first half is a redundant table of contents, mechanically restating theorems and section numbers already covered in the front matter. Furthermore, the ``Further directions'' section is overly myopic. It focuses strictly on internal technical tweaks (e.g., generalizing to a $\sigma$-algebra) rather than establishing macroscopic connections to other major mathematical fields.
    
    \item \textbf{Correction Suggestion:} Completely rewrite the conclusion. Remove the section-by-section summary. Synthesize the core insights and explicitly discuss the framework's broader significance and its macroscopic connections to other established branches of mathematics.
\end{itemize}


\section{Detailed Analysis}

The primary objective of this section is to examine specific cases from the main body of the manuscript and analyze their typographical, stylistic, and logical errors.

\subsection{Analysis of Proof Exposition}

The following excerpt from Section 5 of the manuscript illustrates recurring issues with the author's proof writing style:

\begin{quote}
    \small \textit{Let $X=\{a,b,c\}$ and $\mathcal{A}=\mathcal{P}(X)$. Set
    $$R:=\{a,b\},\qquad I:=\{c\},$$
    so $R\cap I=\varnothing$ and $R\cup I=X\in\mathcal{A}$. Define the map $\Pi_R:X\to R$ by
    $$\Pi_R(a):=b,\qquad \Pi_R(b):=a,\qquad \Pi_R(c):=a.$$
    Then $\Pi_R(X)\subseteq R$, but $\Pi_R|_R\ne\mathrm{id}_R$ (it swaps $a$ and $b$) and $\Pi_R\circ\Pi_R\ne\Pi_R$ (since $\Pi_R(\Pi_R(a))=\Pi_R(b)=a\ne b=\Pi_R(a)$). Hence Axiom I fails.}
\end{quote}



\begin{itemize}
    \item \textcolor{NavyBlue}{\textbf{[Language \& Punctuation] Abuse of Parentheses:}} The author excessively relies on nested parentheses, improperly conflating grammatical asides with mathematical notation. This overloading of sentence structure degrades readability. Parenthetical remarks must not be used as a crutch to compress complex logical steps or obscure omitted details.
    
    \item \textcolor{RustOrange}{\textbf{[Style \& Architecture] Obscure Exposition \& Omission of Steps:}} The proof style assumes an unreasonable level of retained context. By relying heavily on long-distance references, unstated field jargon, and the assumption that omitted steps are ``obvious,'' the exposition becomes hostile to readers seeking to grasp the core ideas. Proofs must be self-contained and explicit, not treated as informal shorthand.
    
    \item \textbf{Correction Suggestion:} Unpack all nested parentheses into separate, grammatically complete sentences. Expand compressed logical transitions and explicitly write out steps without relying on omitted context, ensuring the proof can be strictly followed from beginning to end without ambiguity.
\end{itemize}

\textbf{Use of Parentheses}

With a slightly different perspective, parentheses frequently \textbf{mask omissions and implications} rather than providing explicit details, acting as a shortcut for what the author assumes is obvious. Because this perceived \textbf{"obviousness" is highly subjective to the author's mind}, relying on parentheses to hide steps often creates a severe comprehension gap. Therefore, authors must strictly \textbf{avoid nested parentheses and instead treat the reader as a robot.} Like a precise machine, a reader relies on flawless, continuous logic; therefore, even a single vague assumption or omitted detail can cause a complete breakdown of the entire reading process.

\textbf{Dilemma of Mathematical Proof Writing}

The core dilemma of mathematical proof writing lies in \textbf{the cognitive gap between the immersed expert and the general reader}. A proof may seem perfectly concrete to an expert willing to trace long-distance references, decode field-specific jargon, and rely on intuitive leaps. However, to a non-expert or someone simply seeking the central idea, this compressed style is entirely impenetrable. This conflict stems from the author's tendency to omit details they subjectively deem "repetitive" or "obvious." Ultimately, it is the author's fundamental duty to \textbf{bridge this gap by resisting the urge to overly reduce the text, minimizing omissions, and consistently striving for the simplest, most explicit exposition possible.}

\subsection{Typographical Treatment of Symbols}

The following excerpt outlines the global regularity hypotheses used in the manuscript:

\begin{quote}
\small \textit{
Assume Axiom~I holds. Assume moreover that $\{r\}\in\mathcal A$ and $F_r\in\mathcal A$ for every $r\in R$, and that one of the following global regularity hypotheses holds:
\begin{itemize}[nosep]
\item[\emph{(R--fin)}] $R$ is finite; or
\item[\emph{(R--ctbl)}] $R$ is countable and $\mu$ is $\sigma$--additive on $\mathcal A$ (so $\mathcal A$ is assumed to be a $\sigma$--algebra).
\end{itemize}}
\end{quote}
\begin{itemize}
    \item \textcolor{ForestGreen}{\textbf{[Mathematical Formatting] Inconsistent Math Mode for Set Labels:}} In the excerpt above, as well as in designations like ``R-fib'' used elsewhere, the letter ``R'' explicitly refers to the mathematical set $R$. However, it is typeset in standard Roman text within these labels. In rigorous mathematical typography, any character representing a mathematical object must remain in math mode to ensure visual and logical consistency.
    
    \item \textbf{Correction Suggestion:} Always enclose mathematical variables in math mode, even when they appear within text labels or identifiers. Update all such instances to \textbf{($R$-fin)}, \textbf{($R$-ctbl)}, and \textbf{$R$-fib} (using \verb|$R$-fib|) throughout the manuscript.
\end{itemize}

\textbf{Mathematical Typography in Titles}

Including mathematical notation in titles presents a trade-off: it reduces reading cost (e.g., $SL_2(\mathbb{R})$ is vastly more recognizable than spelling it out) but can hinder database indexing. 

Ultimately, two strict principles must be observed: 
\begin{itemize}
    \item Any notation in a title must be universally recognized without additional explanation,
    \item if a concise textual alternative exists, it should always be preferred over a formula.
\end{itemize}
However, textual substitution must never compromise mathematical accuracy. For instance, 
\begin{align*}
    \infty\text{-categories}
\end{align*}
can be flattened into
\begin{center}
    infinity-categories.
\end{center}
But we can not replace 
\begin{align*}
    (\infty,2)\text{-categories}
\end{align*}
by infinity two categories.
\end{document}
